\documentclass[exercice,a4]{thedocument}%
\usepackage[french]{babel}%
\usepackage{pstricks}%
\usepackage{multido}%
\usepackage{scratch3}%
\usepackage{tabularx}%
\usepackage{array}%
\startdatakeys
    \source{22GENMATAS1}
    \BO{2026}
    \cycle{4}
    \competencesDuSocle{CA,RA}
    \objectifsApprentissage{C4E1ag}
    \motsCles{Scratch,programme,polygone}
\enddatakeys
\begin{document}%

\medskip

\emph{Dans cet exercice. aucune justification n'est attendue}

\smallskip

On souhaite réaliser le logo ci-dessous avec le logiciel Scratch à partir du script incomplet ci- dessous.

\begin{center}
\begin{tabularx}{\linewidth}{|X| X|}\hline
\multicolumn{1}{|c|}{Logo}&\multicolumn{1}{|c|}{Script principal}\\
\begin{center}
\psset{unit=1cm,arrowsize=2pt 3}
\begin{pspicture}(-3.3,-3.3)(3.3,3.3)
\def\drapeau{\psline(0,0)(2.3,0)(3.1,-0.75)(1.95,-0.75)(1.15,0)}
\multido{\n=0+72}{5}{\rput{\n}(0;0){\drapeau}}
\end{pspicture}\end{center}&\begin{scratch}[num blocks,num start = 0]
\blockinit{Quand \greenflag est cliqué}
\blockmove{aller à x: \ovalnum0 y: \ovalnum0}
\blockmove{s'orienter à \ovalnum{90} degrés}
\blockpen{effacer tout}
\blockrepeat{répéter \ovalnum{...} fois}
     {\blockmove{Motif}
     \blockmove{aller à x: ... y: ...}
     \blockmove{tourner \turnright{} de \ovalnum{...} degrés}
     }
\end{scratch}\\ \hline
\end{tabularx}
\begin{tabularx}{\linewidth}{|X|}\hline
On rappelle que l'instruction \begin{scratch}\blockmove{s'orienter à \ovalnum{90} degrés}\end{scratch} consiste à orienter le lutin et le stylo horizontalement vers la droite.\\ \hline
Le bloc \begin{scratch}\blockmove{Motif}\end{scratch} permet de réaliser la figure ci-contre :\quad
\psset{unit=1cm,arrowsize=2pt 3}
\begin{pspicture}(-2,-1.2)(3.3,0.2)
\def\drapeau{\psline(0,0)(2.3,0)(3.1,-0.75)(1.95,-0.75)(1.15,0)}
\rput(0;0){\drapeau}
\rput(-1,-1){Point de départ}
\psline[linewidth=2pt]{->}(-1,-0.7)(0,0)
\end{pspicture}\\ \hline
\end{tabularx}
\end{center}

\begin{enumerate}
\item En mathématiques, comment appelle-t-on la transformation géométrique qui permet de
passer d'un motif du logo au suivant ?
\item Ici, le stylo est orienté horizontalement vers la droite au départ.

Parmi les trois propositions suivantes, quelle est celle qui permet d'obtenir le motif souhaité ?

\begin{center}
\begin{tabularx}{\linewidth}{|X| X| X|}\hline
Proposition 1 &Proposition 2& Proposition 3\\ \hline
{\small \begin{scratch}
\initmoreblocks{définir \namemoreblocks{Motif}}
\blockpen{stylo en position d’écriture}
\blockrepeat{répéter \ovalnum{2} fois}
{\blockmove{avancer de \ovalnum{50}}
\blockmove{tourner \turnright{} de \ovalnum{140} degrés}
\blockmove{avancer de \ovalnum{50}}
\blockmove{tourner \turnright{} de \ovalnum{40} degrés}
}
\blockpen{relever le stylo}
\end{scratch}}&{\small \begin{scratch}
\initmoreblocks{définir \namemoreblocks{Motif}}
\blockpen{stylo en position d’écriture}
\blockmove{avancer de \ovalnum{50}}
\blockrepeat{répéter \ovalnum{2} fois}
{\blockmove{avancer de \ovalnum{50}}
\blockmove{tourner \turnright{} de \ovalnum{40} degrés}
\blockmove{avancer de \ovalnum{50}}
\blockmove{tourner \turnright{} de \ovalnum{140} degrés}
}
\blockpen{relever le stylo}
\end{scratch}}&{\small \begin{scratch}
\initmoreblocks{définir \namemoreblocks{Motif}}
\blockpen{stylo en position d’écriture}
\blockmove{avancer de \ovalnum{50}}
\blockrepeat{répéter \ovalnum{2} fois}
{\blockmove{avancer de \ovalnum{50}}
\blockmove{tourner \turnleft{} de \ovalnum{40} degrés}
\blockmove{avancer de \ovalnum{50}}
\blockmove{tourner \turnleft{} de \ovalnum{140} degrés}
}
\blockpen{relever le stylo}
\end{scratch}}\\ \hline
\end{tabularx}
\end{center}


\item Compléter le script principal en recopiant sur la copie uniquement la boucle « répéter» (c'est-à-dire les instructions 4, 5, 6 et 7).
\item On veut placer l'instruction \begin{scratch} \blockpen{ajouter \ovalnum{10} à la couleur du style}\end{scratch} de façon à changer de couleur à chaque motif.

Sur la copie, indiquer un numéro d'instruction du script principal après laquelle on peut placer cette instruction.
\end{enumerate}

\medskip

\showthesource
\end{document}
